problem:  
Let $(M^n,g)$ be a closed Riemannian manifold satisfying the **harmonic‑curvature** condition $\nabla^*\mathrm{Riem}=0$.  Denote by  
\[
\mathcal{H}(M)=\{\,g\in\mathrm{Met}(M)\mid\nabla^*_g\mathrm{Riem}_g=0\,\}/\mathrm{Diff}(M)
\]
the moduli space of harmonic‑curvature metrics modulo diffeomorphisms.  For an Einstein metric $g_E$, standard theory tells us that its infinitesimal deformations satisfy the linearized Einstein operator, which is elliptic modulo gauge and hence admits a finite‑dimensional Kuranishi model.

**Question:**  
Does there exist a compact, non‑Einstein harmonic‑curvature metric $g_0$ on some manifold $M$ that is **locally rigid** in the sense that every sufficiently small smooth perturbation $g_t=g_0+t\,h+O(t^2)$ satisfying $\nabla^*_{g_t}\mathrm{Riem}_{g_t}=0$ is obtained from $g_0$ by a diffeomorphism and scaling?  
Equivalently, is the linearized operator
\[
D_{g_0}(\nabla^*\mathrm{Riem}) : \Gamma(S^2T^*M)\to \Gamma(T^*M\otimes\Lambda^2T^*M)
\]
elliptic modulo diffeomorphism invariance with **trivial kernel** at some non‑Einstein $g_0$?  
If no such example exists, must every non‑Einstein harmonic‑curvature metric admit a genuine deformation family within $\mathcal{H}(M)$, implying analytic “flexibility’’ analogous to the deformation theory of self‑dual or conformally Einstein metrics?

Why is it a "good" problem:  
This version sharpens and corrects the focus toward **rigidity vs. flexibility** of harmonic‑curvature metrics, without contradicting elliptic PDE expectations.  It asks whether non‑Einstein harmonic‑curvature metrics can ever be *isolated* points in their moduli space—an analytically and geometrically subtle question.  

A positive answer would show that harmonic‑curvature metrics mirror the strong rigidity of symmetric or Einstein structures, while a negative one would reveal new deformation phenomena between curvature‑harmonic and conformally Einstein geometries.  Either outcome would deepen our understanding of the analytic structure of $\nabla^*\mathrm{Riem}=0$, clarifying whether its moduli space behaves more like the rigid Einstein moduli or the flexible self‑dual/conformal ones.  

The problem unites **Riemannian geometry**, **elliptic PDE analysis**, and **deformation theory**, providing a clean, well‑posed, yet genuinely open conjectural direction.  Its appeal lies in simplicity of formulation—“Are non‑Einstein harmonic‑curvature metrics ever rigid?”—and depth of consequence, making it a truly “good’’ research‑level mathematical problem.